\documentclass[12pt,a4paper]{article}

\usepackage[utf8]{inputenc}
\usepackage[T1]{fontenc}
\usepackage{amsmath,amssymb,amsthm,mathtools}
\usepackage[colorlinks=true,linkcolor=blue,citecolor=red,urlcolor=blue]{hyperref}
\usepackage{geometry}
\geometry{margin=2.5cm}
\usepackage{enumitem}
\usepackage{booktabs}
\usepackage[capitalise,noabbrev]{cleveref}

\newtheorem{theorem}{Theorem}[section]
\newtheorem{lemma}[theorem]{Lemma}
\newtheorem{proposition}[theorem]{Proposition}
\newtheorem{corollary}[theorem]{Corollary}
\newtheorem{conjecture}[theorem]{Conjecture}
\theoremstyle{definition}
\newtheorem{definition}[theorem]{Definition}
\theoremstyle{remark}
\newtheorem{remark}[theorem]{Remark}

\newcommand{\Ksym}{K_\sigma^{\mathrm{sym}}}
\newcommand{\PP}{\mathcal{P}}

\title{A Gauge-theoretic Reformulation of the Riemann Hypothesis:\\ Part II: Analytic Obstructions and the Limits of Monotonicity\thanks{This paper is part of the FST-RH program (Zenodo project DOI: 10.5281/zenodo.RH\_DOI\_PLACEHOLDER).}}
\author{Lukas Geiger\\
\small Bernau, Germany\\
\small ORCID: \href{https://orcid.org/0009-0005-7296-1534}{0009-0005-7296-1534}}
\date{March 15, 2026}

\begin{document}
\maketitle

\begin{abstract}
Following the construction of the Prime Hub Graph transfer operator in Part I, this paper rigorously establishes the analytic bridge between the thermodynamic model and the Riemann Hypothesis. We decompose the Li coefficients $\lambda_n$ via the Bombieri--Lagarias framework into archimedean and finite-part contributions, identifying the ``Archimedean Dominance Paradox'': while arithmetic prime-sums alone exhibit instability at the critical point, the completed zeta function is stabilized by the convexity of the Tate-Gamma block. We analyze the prime-sum functionals $A(\sigma), B(\sigma), C(\sigma), D(\sigma)$ and their convergence properties, establishing a boundary-layer regime near $\sigma = 1$ where the covariance of the Gibbs measure dominates the drift. We formulate the remaining gap as a Gauge Positivity problem for the symmetrized arithmetic kernel, and discuss the limitations of this approach.
\end{abstract}

\noindent\textbf{MSC 2020:} 11M26, 11M36, 60F05\\
\textbf{Keywords:} Riemann Hypothesis, Li coefficients, archimedean dominance, gauge positivity, prime zeta function

\tableofcontents

\section{Introduction}

Li's criterion asserts that the Riemann Hypothesis is true if and only if $\lambda_n \ge 0$ for all $n \ge 1$. In Part I, we mapped the zeros of $\zeta(s)$ to the spectrum of a transfer operator. Here, we transition from heuristics to a rigorous analytical framework, transforming the RH into a stability problem of a renormalized arithmetic constraint.

\section{The Renormalized Framework}

We define the objects required for the formal proof. 

\begin{definition}[Li Observable and Normalization]
The arithmetic Gibbs measure is $\mu_\sigma(p) = p^{-\sigma}/P(\sigma)$ with $P(\sigma) = \sum_p p^{-\sigma}$. The $n$-th Li observable is defined as:
\[ X_{n,\sigma}(p) = \sum_{k=1}^{n} \binom{n}{k} \frac{(n-1)!}{(k-1)!} (-1)^k (\log p)^k \operatorname{Li}_{1-k}(p^{-\sigma}). \]
The renormalized constraint is $m_n(\sigma) = \mathbb{E}_{\mu_\sigma}[X_{n,\sigma}] - h_n(\sigma)$.
\end{definition}

\begin{definition}[Archimedean Bulk]
The archimedean subtraction term $h_n(\sigma)$ is defined in closed form from the non-arithmetic part of the completed zeta function:
\[ h_n(\sigma) = \frac{1}{(n-1)!} \frac{\partial^n}{\partial s^n} \left[ s^{n-1} \log \left( \frac{1}{2} \pi^{-s/2} \Gamma(s/2) s \right) \right]_{s=\sigma}. \]
\end{definition}

\section{Preparatory Lemmas}

The Critical Identification Theorem below requires control over the divergences that arise as $\sigma \to 1^+$. We establish two foundational lemmas.

\begin{lemma}[Archimedean Cancellation]\label{lem:arch-cancel}
For each $n \ge 1$, the Gibbs expectation $\mathbb{E}_{\mu_\sigma}[X_{n,\sigma}]$ admits a pole-subtracted representation. For $n=1$:
\[ \mathbb{E}_{\mu_\sigma}[X_{1,\sigma}] = -\frac{A(\sigma)}{P(\sigma)} = -\frac{1}{P(\sigma)} \sum_p (\log p)\, p^{-\sigma}\, f(p,\sigma) - \frac{B(\sigma)}{(\sigma-1)\,P(\sigma)}, \]
where $f(p,\sigma) := \frac{1}{p^\sigma - 1} - \frac{1}{(\sigma-1)\log p}$ satisfies $\lim_{\sigma \to 1^+} f(p,\sigma) = \frac{1}{p-1} - \frac{1}{\log p}$. The first sum converges to a finite limit, while the second term captures the divergence of order $1/(\sigma-1)$.
\end{lemma}
\begin{proof}
The Laurent expansion $\frac{1}{e^u-1} = \frac{1}{u} - \frac{1}{2} + \frac{u}{12} - \cdots$ with $u = \sigma \log p$ yields the stated limit. For large $p$, $|f(p,1)| = |\frac{1}{p-1} - \frac{1}{\log p}| \le C/p$ since $\log p \ll p$. Therefore
\[ \sum_p \frac{\log p}{p}\,|f(p,1)| \le C' \sum_p \frac{\log p}{p^2} < \infty \]
by partial summation against the Chebyshev bound $\pi(t) \le 2t/\log t$.
\end{proof}

\begin{remark}[Gibbs Normalization Gap]\label{rem:gibbs-gap}
Since the numerator sum $S(\sigma) = \sum_p (\log p)\,p^{-\sigma}\,f(p,\sigma)$ converges to a finite limit while $P(\sigma) \to \infty$, the Gibbs-normalized ``finite remainder'' vanishes: $-S(\sigma)/P(\sigma) \to 0$. This means $\mathbb{E}_{\mu_\sigma}[X_{1,\sigma}]/P(\sigma)$ loses the information encoded in $\lambda_1^{\mathrm{fin}} = -\log 2$. The Li coefficients arise from $\frac{d}{ds}\log\xi(s)\big|_{s=1}$, which involves $\zeta'(s)/\zeta(s) + 1/(s-1) \to \gamma$ as $s \to 1^+$ --- a cancellation that occurs \emph{before} Gibbs normalization, not after. This observation motivates the direct approach via $\log\xi$ in \Cref{thm:B21}.
\end{remark}

\begin{lemma}[Finite-Part Convergence]\label{lem:finite-part}
The following prime sums converge absolutely at $\sigma = 1$ and are right-continuous:
\begin{align*}
A(\sigma) &:= \sum_p \frac{\log p}{p^\sigma(p^\sigma - 1)}, \quad
C(\sigma) := \sum_p \frac{(\log p)^2}{p^\sigma(p^\sigma - 1)}, \quad
D(\sigma) := \sum_p \frac{(\log p)^2}{(p^\sigma - 1)^2}.
\end{align*}
\end{lemma}
\begin{proof}
For $p \ge 2$ and $\sigma \ge 1$: $p^\sigma - 1 \ge p/2$, giving $\frac{1}{p^\sigma(p^\sigma-1)} \le \frac{2}{p^2}$ and $\frac{1}{(p^\sigma-1)^2} \le \frac{4}{p^2}$. By partial summation against $\pi(t) \le 2t/\log t$:
\[ \sum_p \frac{(\log p)^j}{p^2} \le C_j \int_2^\infty \frac{(\log t)^{j-1}}{t^2}\,dt < \infty \quad (j=1,2). \]
Right-continuity follows from dominated convergence, since each summand is continuous in $\sigma$ and the dominating series $\sum_p (\log p)^2/p^2$ is $\sigma$-independent.
\end{proof}

\section{The Li Coefficients via the Completed Zeta Function}

The connection between the thermodynamic framework and the Riemann Hypothesis is mediated by the completed xi function $\xi(s) = \tfrac{1}{2}s(s-1)\pi^{-s/2}\Gamma(s/2)\zeta(s)$.

\begin{theorem}[Bombieri--Lagarias Representation] \label{thm:B21}
For any integer $n \ge 1$, the $n$-th Li coefficient admits the representation:
\[ \lambda_n = \frac{1}{(n-1)!} \frac{\partial^n}{\partial s^n} \Big[ s^{n-1} \log \xi(s) \Big]_{s=1}. \]
Since $\xi(s)$ is entire of order $1$ with $\xi(1) = 1/2 \ne 0$, the function $\log\xi(s)$ is holomorphic in a neighborhood of $s=1$, and the above derivative is well-defined.
\end{theorem}

\begin{proof}
This is the content of \cite{BombieriLagarias1999}, Theorem 1. The key identity is $\lambda_n = \sum_\rho [1 - (1-1/\rho)^n]$, which follows from the Hadamard product representation of $\xi$.
\end{proof}

\begin{proposition}[Archimedean--Finite Decomposition]\label{prop:arch-fin-decomp}
Each Li coefficient decomposes as $\lambda_n = \lambda_n^{(\infty)} + \lambda_n^{(\mathrm{fin})}$, where:
\begin{align}
\lambda_n^{(\infty)} &= \frac{1}{(n-1)!} \frac{\partial^n}{\partial s^n} \Big[ s^{n-1} \log\!\big(\tfrac{1}{2}s\,\pi^{-s/2}\,\Gamma(s/2)\big) \Big]_{s=1}, \label{eq:arch-part} \\
\lambda_n^{(\mathrm{fin})} &= \frac{1}{(n-1)!} \frac{\partial^n}{\partial s^n} \Big[ s^{n-1} \log g(s) \Big]_{s=1}, \quad g(s) := (s-1)\zeta(s). \label{eq:fin-part}
\end{align}
The function $g(s)$ is holomorphic at $s=1$ with $g(1)=1$, so $\log g$ is holomorphic and all derivatives in~\eqref{eq:fin-part} are finite. The ``pole contribution'' from $\log(s-1)$ cancels between $\log\xi$ and $\log(\frac{1}{2}s\cdot\pi^{-s/2}\Gamma(s/2)) + \log g(s)$, with the residual cross-term absorbed into $\lambda_n^{(\infty)}$.
\end{proposition}

\begin{proof}
Since $\xi(s) = \frac{1}{2}s(s-1)\pi^{-s/2}\Gamma(s/2)\zeta(s)$ and $(s-1)\zeta(s) = g(s)$, we have $\log\xi(s) = \log(\frac{1}{2}s\,\pi^{-s/2}\Gamma(s/2)) + \log g(s)$. Both functions on the right are holomorphic at $s=1$, and the decomposition follows by linearity of differentiation.
\end{proof}

\begin{remark}
For $n=1$: $\lambda_1^{(\infty)} = 1 - \frac{\log\pi}{2} + \frac{\psi(1/2)}{2} = 1 - \frac{\log\pi}{2} - \frac{\gamma}{2} - \log 2 \approx -0.554$ and $\lambda_1^{(\mathrm{fin})} = g'(1)/g(1) = \gamma \approx 0.577$, giving $\lambda_1 = 1 + \gamma/2 - \log(4\pi)/2 \approx 0.023 > 0$.
\end{remark}

\begin{remark}[Asymptotic Behavior of the Decomposition]\label{rem:asymp-decomp}
Numerical computation for $n = 1, \ldots, 18$ reveals that $\lambda_n^{(\mathrm{fin})} > 0$ throughout, with a peak at $n \approx 6$ ($\lambda_6^{(\mathrm{fin})} \approx 1.488$) and subsequent decay $\lambda_n^{(\mathrm{fin})} \approx 10.4\,n^{-0.89}$ for large~$n$. Meanwhile, $\lambda_n^{(\infty)}$ crosses from negative to positive near $n = 8$ and grows as $\lambda_n^{(\infty)} \sim \frac{n}{2}\log n$. This creates a three-regime structure:
\begin{enumerate}[label=(\alph*)]
  \item \textbf{Small $n$ ($n \lesssim 8$):} $\lambda_n^{(\infty)} < 0$, compensated by $\lambda_n^{(\mathrm{fin})} > 0$. Positivity is a delicate cancellation.
  \item \textbf{Transition ($n \approx 8$):} $\lambda_n^{(\infty)}$ crosses zero; both parts contribute positively.
  \item \textbf{Large $n$ ($n \gg 10$):} $\lambda_n^{(\infty)}$ dominates; positivity is automatic.
\end{enumerate}
The ``hard'' instances of Li's criterion are therefore concentrated in regime~(a). Whether $\lambda_n^{(\mathrm{fin})} > 0$ for all $n$ is an open question (see Part~IV, OP2). The decay of $\lambda_n^{(\mathrm{fin})}$ is governed by the radius of convergence $R = 3$ of $\log g(s)$ at $s=1$ (determined by the trivial zero of $\zeta$ at $s=-2$).
\end{remark}

\begin{remark}[Connection to the Gibbs Framework]\label{rem:gibbs-connection}
The Gibbs expectation $\mathbb{E}_{\mu_\sigma}[X_{1,\sigma}] = -A(\sigma)/P(\sigma)$ from Definition~2.1 captures the prime-sum structure but is \emph{not} a regularization of $\zeta'(\sigma)/\zeta(\sigma)$. As shown in \Cref{rem:gibbs-gap}, the $1/P(\sigma)$-normalization annihilates the finite-part contribution. The direct decomposition of $\log\xi$ via \Cref{prop:arch-fin-decomp} avoids this issue entirely.
\end{remark}

\section{The Archimedean Dominance Paradox}

This section identifies the fundamental stability mechanism of the Riemann Hypothesis.

\begin{proposition}[Stability Decomposition]
The global curvature $\mathcal{M}(\sigma) = \partial_s^2 \log \xi(s)$ decomposes into an arithmetic and an archimedean part:
\[ \mathcal{M}(\sigma) = \mathcal{M}_{\text{arith}}(\sigma) + \mathcal{M}_{\text{arch}}(\sigma). \]
Where:
\begin{enumerate}
    \item \textbf{Arithmetic Instability:} $\mathcal{M}_{\text{arith}}(1) = \sum_{\rho} \Re(1/\rho^2)$. Since the first zero is at $\gamma_1 \approx 14.13$, this sum is negative, showing the primes alone are unstable.
    \item \textbf{Archimedean Stability:} $\mathcal{M}_{\text{arch}}(\sigma)$ is explicitly positive and convex for $\sigma > 1/2$.
\end{enumerate}
\textbf{Conclusion:} The Riemann Hypothesis is equivalent to the statement that the archimedean convexity $\mathcal{M}_{\text{arch}}$ globally dominates the arithmetic concavity $\mathcal{M}_{\text{arith}}$ for all $\sigma > 1/2$. This dominance is a global stability requirement that does not follow locally from the functional equation alone.
\end{proposition}

\section{Local Analysis of Prime-Sum Functionals}

While the $\log\xi$ approach of the previous section provides the definitive representation, the prime-sum functionals from Definition~2.1 have independent analytical interest. We analyze the target functional $F(\sigma) = A(\sigma)B(\sigma) - P(\sigma)[C(\sigma)+D(\sigma)]$ which arises from differentiating the Gibbs variance.

\begin{proposition}[Boundary Layer for $F(\sigma)$]\label{prop:boundary-layer}
The target functional $F(\sigma) > 0$ for $\sigma$ in a right-neighborhood $(1, \sigma_0)$ with $\sigma_0 \approx 1.14$, and $F(\sigma) < 0$ for $\sigma > \sigma_0$.
\end{proposition}
\begin{proof}
By \Cref{lem:finite-part}, the sums $A(\sigma), C(\sigma), D(\sigma)$ converge at $\sigma=1$ and are right-continuous. Near $\sigma = 1 + \varepsilon$: the product $A \cdot B$ scales as $\sim 1/\varepsilon$ (since $B(\sigma) \sim \log(1/\varepsilon)/\varepsilon$), while $P \cdot (C+D)$ scales as $\sim \log(1/\varepsilon)$. Thus $F(\sigma) > 0$ for sufficiently small $\varepsilon$. Numerical computation with $N = 10{,}000$ primes (including rigorous tail bounds via the Rosser--Schoenfeld estimate) confirms $F(\sigma) > 0$ for $\sigma \in (1, 1.14)$ and $F(\sigma) < 0$ for all tested $\sigma > 1.14$, with the minimum $F \approx -0.184$ at $\sigma \approx 1.28$.
\end{proof}

\begin{remark}
This sign change of $F(\sigma)$ is a fundamental feature, not a numerical artifact. It reflects the fact that the Gibbs-normalized prime-sum formulation captures only a local approximation to the Li coefficients (cf.\ \Cref{rem:gibbs-gap}), and cannot serve as a global proxy for the $\log\xi$ derivative.
\end{remark}

\section{Structural Results: Non-Identification, Variance Bounds, and Uncertainty}

We establish three unconditional results that constrain the relationship between the thermodynamic and spectral approaches.

\begin{proposition}[Non-Identification]\label{prop:non-ident}
The excess free energy $I_n := \lim_{\sigma \to 1^+} \Phi_n(\sigma)$ and the Li coefficient $\lambda_n$ are generically \emph{not} equal:
\[ I_n \neq \lambda_n. \]
This is because $I_n$ is a \emph{tilt cost} (KL divergence under the Gibbs measure $\mu_\sigma$), while $\lambda_n$ is a \emph{trace functional} (unweighted spectral invariant of $\log\xi$). The Gibbs expectation $\omega_\sigma(H) = P(\sigma)^{-1}\sum_p p^{-\sigma}H(p)$ and the trace $\sum_p H(p)$ are structurally distinct.
\end{proposition}

\begin{proposition}[Variance Lower Bound for $I_1$]\label{prop:variance-bound}
For $\sigma > 1$, the excess free energy satisfies:
\[ \Phi_1(\sigma) \geq \tfrac{1}{2}\,\mathrm{Var}_{\mu_\sigma}(H_{1,\sigma}), \]
where $H_{1,\sigma}(p) = -\log(1 - p^{-\sigma})$. In particular, $I_1 \geq \tfrac{1}{2}\,V_1$ where $V_1 := \lim_{\sigma \to 1^+}\mathrm{Var}_{\mu_\sigma}(H_{1,\sigma}) > 0$.
\end{proposition}
\begin{proof}
The cumulant generating function satisfies $\log\,\mathbb{E}[e^X] \geq \mathbb{E}[X] + \frac{1}{2}\mathrm{Var}(X)$ (from the convexity of $e^x$ via Jensen's inequality). Applied to $\Phi_1 = \log\,\mathbb{E}_{\mu}[e^{H_1}] - \mathbb{E}_{\mu}[H_1]$, this gives $\Phi_1 \geq \frac{1}{2}\mathrm{Var}(H_1)$. The limit $V_1$ exists by dominated convergence (partial summation against the PNT). Strict positivity $V_1 > 0$ follows because $H_{1,\sigma}(p) = -\log(1-p^{-\sigma})$ is strictly decreasing in $p$, hence non-constant. Numerically, $V_1 \approx 0.068$, giving $I_1 \geq 0.034 > \lambda_1 \approx 0.023$.
\end{proof}

\begin{theorem}[Arithmetic Uncertainty Relation]\label{thm:uncertainty}
For any normalized $\psi \in \ell^2(\mathcal{P})$, define the prime-domain spread $\Delta_D^2 := \sum_p |a_p|^2(\log p - \bar{x})^2$ and the zero-domain spread $\Delta_t^2 := \int |\hat\psi(t)|^2\,(t - \bar{t})^2\,dt$, where $\hat\psi(t) = \sum_p a_p\,e^{-it\log p}$ is the Fourier-type transform. Then:
\begin{equation}\label{eq:uncertainty}
\Delta_D \cdot \Delta_t \;\geq\; \frac{1}{2}.
\end{equation}
\end{theorem}
\begin{proof}
The substitution $x_p := \log p$ maps $\psi$ to a measure on $\mathbb{R}$ supported at $\{\log p\}$, with $\hat\psi$ as its Fourier transform. The classical Heisenberg--Kennard--Weyl inequality gives $\mathrm{Var}_x(\psi)\cdot\mathrm{Var}_t(\hat\psi) \geq 1/4$, i.e., $\Delta_D \cdot \Delta_t \geq 1/2$.
\end{proof}

\begin{remark}[Interpretation for the Riemann zeros]
The uncertainty relation constrains the joint localization of primes and zeros: sharp localization of zeros (narrow spectral lines) requires many primes to contribute coherently. A zero off the critical line would require a ``localized phase conspiracy'' among finitely many primes, which the arithmetic incommensurability of $\{\log p\}$ makes heuristically implausible. This is a structural consistency check, not a proof.
\end{remark}

\begin{proposition}[Strict Convexity of the Arithmetic Free Energy]\label{prop:strict-convex}
The function $F(\sigma) := -\log\zeta(\sigma)$ is strictly convex on $(1, \infty)$:
\begin{equation}
F''(\sigma) = \sum_p \frac{(\log p)^2\,p^{-\sigma}}{(1-p^{-\sigma})^2} > 0.
\end{equation}
This means the ``arithmetic specific heat'' is everywhere positive: the prime distribution has positive susceptibility at every ``temperature'' $1/\sigma$.
\end{proposition}
\begin{proof}
Direct computation from the Euler product: $F'(\sigma) = -\zeta'(\sigma)/\zeta(\sigma) = \sum_p \log(p)\,p^{-\sigma}/(1-p^{-\sigma})$, and the second derivative is manifestly positive as a sum of positive terms.
\end{proof}

\section{Axiomatic Framework and Route Analysis}\label{sec:axiomatics}

The thermodynamic approach to the Riemann Hypothesis, developed in the companion paper~\cite{GeigerFSTRH2026}, rests on an axiom system A0--A9. We present an honest status assessment, prove that A7 is now resolved, and analyze the consequences for the two proof routes.

\subsection{Axiom Status}

\begin{center}
\renewcommand{\arraystretch}{1.25}
\begin{tabular}{@{}lll@{}}
\toprule
\textbf{Axiom} & \textbf{Status} & \textbf{Content} \\
\midrule
A0 & \textsc{verified} & Hilbert space $\ell^2(\PP)$, scale operator $D = \mathrm{diag}(\log p)$ \\
A1 & \textsc{verified} & Unitary phase flow $U_t = e^{-itD}$ (Stone's theorem) \\
A2 & \textsc{verified} & Amplitude operator $A = \mathrm{diag}(p^{-1/2})$, compact \\
A3 & \textsc{verified} & Transfer operator $T_t = AU_t$, time-reversal $T_t^* = T_{-t}$ \\
A4 & \textsc{verified} & $C^*$-algebra $\mathfrak{A} = C^*(A, \{U_t\})$, state space standard \\
A5 & \textsc{verified} & Gibbs reference state $\omega_\sigma$, invariance from diagonality \\
A6 & \textsc{verified} & Relative entropy = KL divergence ($\mathfrak{A}$ commutative) \\
\midrule
A7 & \textbf{\textsc{proven}} & OS reflection positivity (\Cref{prop:A7-proof} below) \\
A8 & \textbf{\textsc{broken}} & Original formulation inconsistent (\Cref{rem:A8-inconsistency}) \\
A8$'$ & \textbf{\textsc{open}} & Reformulated: nuclear Fredholm det for $\xi$ (\Cref{rem:A8-reformulated}) \\
\midrule
A9(i--ii) & \textsc{verified} & Legendre transform $\Phi_n(\sigma) \ge 0$; critical limit $I_n$ exists \\
A9(iii) & \textbf{\textsc{broken}} & $I_n \neq \lambda_n$; convexity trap (\Cref{rem:A9-broken}) \\
\bottomrule
\end{tabular}
\end{center}

\subsection{Proof of A7: OS Reflection Positivity}

\begin{proposition}[OS Reflection Positivity]\label{prop:A7-proof}
The Osterwalder--Schrader form
\[
  \langle f, g \rangle_\omega
  = \int_0^\infty \!\!\int_0^\infty
    \overline{f(t)}\,g(u)\; K(t,u)\,dt\,du,
  \qquad
  K(t,u) = \sum_{p \in \PP} w_p\, e^{-i(u-t)\log p},
\]
with weights $w_p = p^{-(\sigma+1)}/P(\sigma) > 0$, is positive definite on $L^2_+(\mathbb{R})$. The OS Hilbert space $\mathcal{H}_{\mathrm{OS}}$ carries a contractive semigroup $e^{-\tau H}$ with $H = \mathrm{diag}(\log p)$.
\end{proposition}

\begin{proof}
\emph{Step 1 (Bochner).}
Define the positive measure $\mu := \sum_{p} w_p\,\delta_{\log p}$ on $\mathbb{R}$. Then $K(t,u) = \hat{\mu}(u-t)$, the Fourier transform of $\mu$. By Bochner's theorem, $K$ is positive definite on $\mathbb{R}$: for any $f \in L^2(\mathbb{R})$,
\[
  \int\!\!\int \overline{f(t)}\,f(u)\,K(t,u)\,dt\,du
  = \int |\hat{f}(\xi)|^2\,d\mu(\xi)
  = \sum_p w_p\,|\hat{f}(\log p)|^2 \ge 0.
\]

\emph{Step 2 (Isometry).}
The map $\Phi: L^2_+(\mathbb{R}) \to \ell^2(\PP, w)$ defined by $(\Phi f)_p := \hat{f}(\log p)$ satisfies $\langle f, g\rangle_\omega = \langle \Phi f, \Phi g\rangle_{\ell^2(w)}$. The OS Hilbert space $\mathcal{H}_{\mathrm{OS}} := \overline{\Phi(L^2_+)}$ is isometrically embedded in $\ell^2(\PP, w)$.

\emph{Step 3 (Contractive semigroup).}
Define $H := \mathrm{diag}(\log p)$ on $\ell^2(\PP, w)$. The semigroup $e^{-\tau H} = \mathrm{diag}(p^{-\tau})$ acts as a contraction for $\tau \ge 0$ since $\log p > 0$ for all primes $p$, so $\|e^{-\tau H}\| = \sup_p p^{-\tau} = 2^{-\tau} < 1$. The semigroup leaves $\mathcal{H}_{\mathrm{OS}}$ invariant and is strongly continuous.
\end{proof}

\subsection{Obstruction at A9(iii)}

\begin{remark}[Why A9(iii) is broken]\label{rem:A9-broken}
Two independent obstructions kill the identification $I_n = \lambda_n$ required by A9(iii):

\begin{enumerate}[label=(\alph*)]
\item \textbf{K1 (Non-Identification).} The excess free energy $I_n = \lim_{\sigma \to 1^+} \Phi_n(\sigma)$ is a tilt cost under the Gibbs measure $\mu_\sigma$, while $\lambda_n$ is an unweighted spectral trace of $\log\xi$. As shown in \Cref{prop:non-ident}, $I_n \neq \lambda_n$ generically. The Gibbs normalization $1/P(\sigma)$ annihilates the finite-part information (\Cref{rem:gibbs-gap}).

\item \textbf{Convexity Trap.} Even granting $I_n \ge 0$ (which holds unconditionally from the KL interpretation), this does \emph{not} imply $\lambda_n \ge 0$. By the Bregman inequality (\Cref{rem:convexity-trap}), $I_n \ge 0$ is a universal identity for any strictly convex generating function, independent of $\mathrm{sign}(\lambda_n)$.
\end{enumerate}

\noindent
Taken together, Route~2 (proving RH via A9) is dead: A9(iii) cannot be repaired without an independent proof of Li positivity.
\end{remark}

\subsection{Route Analysis}

\begin{remark}[Two routes to the Riemann Hypothesis]\label{rem:route-analysis}
The axiom system admits two logical routes to RH:

\begin{description}[font=\bfseries,leftmargin=2em]
\item[Route 1 (A7 $+$ A8$'$):]
A7 (OS reflection positivity) is now proven (\Cref{prop:A7-proof}). The original A8 is inconsistent (\Cref{rem:A8-inconsistency}) and must be replaced by A8$'$ (\Cref{rem:A8-reformulated}): the construction of a nuclear operator $\mathcal{L}_s$ with $\xi(s) = C \cdot \det(I - \mathcal{L}_s)$.

\item[Route 2 (A9):]
Dead. A9(i--ii) hold, but A9(iii) is broken by the K1 non-identification and the convexity trap (\Cref{rem:A9-broken}).
\end{description}

\noindent
\textbf{Conclusion.} The Riemann Hypothesis is now localized to A8$'$.
\end{remark}

\subsection{K4: Inconsistency of the Original A8}

\begin{remark}[A8 is inconsistent with Hardy's theorem]\label{rem:A8-inconsistency}
The original A8 posits a trace-class positive operator $\mathcal{K}_0 \ge 0$ such that
$\xi(1/2+it) = C \cdot \det(I + V_t\,\mathcal{K}_0\,V_t^*)$
with $V_t$ unitary. Since $\mathcal{K}_t := V_t\,\mathcal{K}_0\,V_t^* \ge 0$ and is trace class, the Fredholm determinant satisfies
\[
  \det(I + \mathcal{K}_t) = \prod_{j} (1 + \lambda_j(\mathcal{K}_t)) \ge 1 > 0
\]
for all real $t$ (each eigenvalue $\lambda_j \ge 0$, so each factor $\ge 1$). This would imply $\xi(1/2+it) > 0$ for all $t \in \mathbb{R}$.

However, Hardy (1914) proved that $\xi$ has \emph{infinitely many} zeros on the critical line: there exist real $\gamma_n$ with $\xi(1/2 + i\gamma_n) = 0$. Therefore A8 as stated \textbf{contradicts} Hardy's theorem and cannot hold for any choice of $\mathcal{K}_0$, $\mathcal{V}$, $V_t$.
\end{remark}

\subsection{Reformulated A8$'$: Nuclear Fredholm Determinant}

\begin{remark}[Corrected axiom A8$'$]\label{rem:A8-reformulated}
The correct operator-theoretic formulation of the Riemann Hypothesis via Fredholm determinants must use $\det(I - \mathcal{L}_s)$ (which \emph{can} vanish) rather than $\det(I + \mathcal{K}_s)$ with $\mathcal{K}_s \ge 0$ (which cannot).

\medskip\noindent
\textbf{A8$'$ (Nuclear transfer operator for $\xi$).}
There exists a separable Hilbert space $\mathcal{V}$ and a nuclear (trace-class) operator-valued holomorphic function $s \mapsto \mathcal{L}_s$ on $\mathcal{V}$ such that:
\begin{enumerate}[label=(\roman*),nosep]
\item \textbf{Determinant identity:} $\xi(s)/\xi(0) = \det(I - \mathcal{L}_s)$.
\item \textbf{Spectral control:} For $\re(s) > 1/2$ (excluding $s=1$), the spectral radius $\rho(\mathcal{L}_s) < 1$, so $\det(I - \mathcal{L}_s) \ne 0$.
\item \textbf{Zeros on the critical line:} At $s = 1/2 + i\gamma_n$, the operator $\mathcal{L}_s$ acquires an eigenvalue equal to $1$, giving $\det = 0$.
\end{enumerate}
This is the structure of the Selberg zeta function $Z(s) = \det(I - \mathcal{L}_s^{\mathrm{Selberg}})$ for hyperbolic surfaces \cite{Ruelle1976,Mayer1991}, where the transfer operator is nuclear for $\re(s) > 1/2$. For the Riemann zeta, the analogous construction faces a fundamental trace-class obstruction.
\end{remark}

\begin{remark}[The trace-class obstruction for A8$'$]\label{rem:trace-class-obstruction}
The Euler product gives $1/\zeta(s) = \det(I - L_s)$ with $L_s = \mathrm{diag}(p^{-s})$ on $\ell^2(\PP)$. However, $L_s$ is trace class only for $\re(s) > 1$:
\[
  \|L_s\|_1 = \sum_p p^{-\re(s)}
  \begin{cases}
    < \infty & \text{if } \re(s) > 1, \\
    = +\infty & \text{if } \re(s) \le 1.
  \end{cases}
\]
At $s = 1/2 + it$ (the critical line), $\|L_s\|_1 = \sum_p p^{-1/2} = +\infty$, so the Fredholm determinant $\det(I - L_s)$ is not defined.

The completed function $\xi(s)$ incorporates the archimedean factor $\frac{1}{2}s(s-1)\pi^{-s/2}\Gamma(s/2)$, which provides analytic regularization. In the Selberg case, the corresponding regularization comes from the exponential decay of geodesic lengths, making the transfer operator nuclear for $\re(s) > 1/2$. For the Riemann zeta, the primes are ``too dense'' ($\pi(x) \sim x/\log x$) compared to geodesic lengths ($\#\{\gamma : l_\gamma \le T\} \sim e^T/T$), and this density gap is the fundamental reason why the transfer operator approach works for Selberg but not directly for Riemann.

Overcoming this obstruction likely requires incorporating the archimedean place into the operator structure itself (cf.\ Deninger \cite{Deninger2002}, Connes \cite{Connes1999}).
\end{remark}

\section{Generating Function and Path D2}\label{sec:path-d2}

The Bombieri--Lagarias representation suggests a generating function approach to Li positivity.

\begin{proposition}[Generating Function Identity]\label{prop:gen-func}
Define $c_n := \lambda_n / n$. Then:
\begin{equation}\label{eq:gen-func}
\sum_{n=1}^{\infty} c_n\, z^n = \log\xi\!\left(\frac{1}{1-z}\right) - \log\xi(1), \qquad |z| < 1.
\end{equation}
Li's criterion $\lambda_n \geq 0$ for all $n$ is equivalent to $c_n \geq 0$ for all $n$.
\end{proposition}
\begin{proof}
The Li functional applied to an analytic function $f$ holomorphic at $s=1$ extracts the Taylor coefficients of $f(1/(1-z))$ around $z=0$. Applied to $\log\xi$, this gives the identity.
\end{proof}

\begin{remark}[Path D2: Stieltjes Realization --- Obstruction]
If the generating function $F(z) = \sum c_n z^n$ admitted a representation
$F(z) = \int_0^1 \frac{z}{1-wz}\,d\mu(w)$
for some positive measure $\mu$ on $[0,1]$, then $c_n = \int w^n\,d\mu(w) \geq 0$ automatically (proving RH), and $\{c_n\}$ would be \emph{completely monotone}: $(-1)^k \Delta^k c_n \geq 0$ for all $k, n$. However, numerical computation shows that the sequence $c_n = \lambda_n/n$ is \emph{strictly increasing} (see \Cref{sec:numerics}): $c_1 \approx 0.023$, $c_{13} \approx 0.294$, with $c_n \sim \frac{1}{2}\log n$ asymptotically. This violates complete monotonicity already at $k=1$ (the first forward difference $\Delta c_n = c_{n+1} - c_n > 0$, so $(-1)^1 \Delta c_n < 0$). Thus, the ``Stieltjes realization'' approach with a measure on $[0,1]$ is \emph{ruled out}.

Under RH, the correct spectral measure is supported on the \emph{unit circle} $|w_\rho| = 1$ (not the interval $[0,1]$), and $c_n = \frac{1}{n}\sum_\rho [1-(1-1/\rho)^n]$ involves cancellation of oscillatory terms. The positivity of $c_n$ is therefore an essentially \emph{number-theoretic} property of the zero distribution, not a consequence of general positivity theory.
\end{remark}

\section{Spectral Structure of the Arithmetic Kernel}

We discuss the spectral structure of the symmetrized arithmetic kernel $\Ksym$ truncated to $N$ primes.

\begin{proposition}[Spectral Census]\label{prop:spectral-census}
For the symmetrized arithmetic kernel $\Ksym$ truncated to $N \ge 300$ primes and $\sigma \in (1, 5)$, numerical computation shows that the negative inertia index is exactly $2$, independent of $N$ and $\sigma$. In the boundary layer $\sigma \in (1, \sigma_0)$, the total sum $\mathbf{1}^T \Ksym\, \mathbf{1} > 0$ despite these two negative eigenvalues.
\end{proposition}

\begin{remark}[Limitations of the Gauge Approach]
While the rank-2 negative inertia suggests a compact gauge correction might restore positive semi-definiteness for each fixed $\sigma$, the sign change $F(\sigma_0) = 0$ identified in \Cref{prop:boundary-layer} means that the sum $\mathbf{1}^T \Ksym\, \mathbf{1}$ itself becomes negative for $\sigma > \sigma_0$. This rules out a direct reduction of the global RH to a simple gauge positivity condition on $\Ksym$. The spectral census remains a valid numerical observation, but its connection to the Li coefficients requires the $\log\xi$ approach of \Cref{thm:B21}, not the Gibbs-normalized prime-sum formulation.
\end{remark}

\begin{remark}[Excess Free Energy and the Convexity Trap]\label{rem:convexity-trap}
The excess free energy $\Phi_n(\sigma) := \log\,\mathbb{E}_{\mu_\sigma}[e^{X_\sigma}] - \mathbb{E}_{\mu_\sigma}[X_\sigma]$ satisfies $I_n := \lim_{\sigma \to 1^+} \Phi_n(\sigma) \ge 0$, since $\Phi_n(\sigma) = \mathrm{KL}(\nu_{n,\sigma}\|\mu_\sigma)$ for the tilted measure $\nu_{n,\sigma} \propto \mu_\sigma \cdot e^{H_{n,\sigma}}$ \cite{DonskerVaradhan1975}. Existence of $I_n$ follows from a cancellation lemma: while the affine part of $X_\sigma$ diverges as $\sigma \to 1^+$, cumulants of order $r \ge 2$ are affine-invariant, so $\kappa_r(X_\sigma) = \kappa_r(Q_\sigma)$ where $Q_\sigma$ has bounded coefficients. The series $I_n = \sum_{r \ge 2} \kappa_r^*/r!$ converges by dominated convergence.

However, $I_n \ge 0$ does \emph{not} imply $\lambda_n \ge 0$. By the Bregman inequality for any strictly convex $F$: $I_n = F(\Omega_n) - F'(0)\Omega_n \ge 0$ holds universally, regardless of $\mathrm{sign}(F'(0)) = \mathrm{sign}(\lambda_n)$. This ``convexity trap'' shows that no purely thermodynamic convexity argument can force Li positivity --- the gauge-theoretic approach of Proposition~\ref{prop:gauge-pos} is necessary.
\end{remark}

The gap between the Gibbs-normalized prime-sum analysis (which captures local structure near $\sigma = 1$) and the global Li positivity (which requires control of $\log\xi$ at $s=1$) remains the central open problem. The Archimedean--Finite Decomposition of \Cref{prop:arch-fin-decomp} provides the correct framework for addressing it.

\appendix
\section{Numerical Details}\label{sec:numerics}
Numerical tests using $N=10{,}000$ primes with $50$-digit mpmath precision, supplemented by Rosser--Schoenfeld tail bounds, yield: $F(\sigma) > 0$ for $\sigma \in (1.005, 1.14)$ with the transition at $\sigma_0 \approx 1.14$, and $F(\sigma) < 0$ for $\sigma \in (1.14, 20)$. The minimum occurs at $\sigma \approx 1.28$ with $F \approx -0.184$.

The Li coefficients, computed directly from $\log\xi$ via high-precision numerical differentiation at $s = 1 + 10^{-10}$ (to avoid the pole), together with their archimedean--finite decomposition:

\medskip
\begin{center}
\begin{tabular}{r|r|r|r|r}
$n$ & $\lambda_n$ & $\lambda_n^{(\infty)}$ & $\lambda_n^{(\mathrm{fin})}$ & $c_n = \lambda_n/n$ \\
\hline
1 & $0.02310$ & $-0.55412$ & $0.57722$ & $0.02310$ \\
2 & $0.09235$ & $-0.87454$ & $0.96689$ & $0.04617$ \\
3 & $0.20764$ & $-1.01306$ & $1.22070$ & $0.06921$ \\
5 & $0.57554$ & $-0.88273$ & $1.45827$ & $0.11511$ \\
8 & $1.46576$ & $+0.02090$ & $1.44486$ & $0.18322$ \\
10 & $2.27934$ & $+0.95554$ & $1.32380$ & $0.22793$ \\
13 & $3.81724$ & $+2.72452$ & $1.09272$ & $0.29363$ \\
\end{tabular}
\end{center}

\medskip\noindent
Key observations: (i)~All $\lambda_n > 0$ for $n = 1, \ldots, 13$, consistent with Li's criterion. (ii)~The archimedean part $\lambda_n^{(\infty)}$ crosses from negative to positive near $n = 8$, reflecting the transition from Gamma-dominated to polynomial-dominated growth. (iii)~The finite part $\lambda_n^{(\mathrm{fin})}$ is always positive and slowly decreasing from $\gamma \approx 0.577$. (iv)~The sequence $c_n = \lambda_n/n$ is \emph{strictly increasing}, ruling out complete monotonicity and the Stieltjes realization approach (cf.\ \Cref{sec:path-d2}). (v)~The growth ratio $\lambda_n/(n\log n/2)$ increases toward $1$, consistent with the asymptotic $\lambda_n \sim \frac{n}{2}\log n$ under RH.

Scripts are available in the companion repository.

\bibliographystyle{plain}
\bibliography{fst-rh-references}

\end{document}
